\begin{table}[!htbp]\centering
\caption{受限局部水平族的参数支持范围}\label{tab:round2calibration}
\small\setlength{\tabcolsep}{4pt}
\begin{tabularx}{\linewidth}{L r r r r r}\toprule
参数 & 最优 & 支持下限 & 支持上限 & 删一下限 & 删一上限 \\
\midrule
qc & 0.000000 & 0.000000 & 0.000736 & 0.000000 & 0.000000 \\
qd & 0.005888 & 0.002944 & 0.011776 & 0.002944 & 0.005888 \\
r & 0.000736 & 0.000147 & 0.001472 & 0.000736 & 0.000736 \\
rho & 0.500000 & -0.500000 & 0.500000 & 0.500000 & 0.500000 \\
m & 0.000000 & 0.000000 & 0.000000 & 0.000000 & 0.000000 \\
q\_r\_ratio & 4.000000 & 1.000000 & 45.000000 & 2.000000 & 4.000000 \\
\bottomrule\end{tabularx}
\par\vspace{3pt}\begin{minipage}{\linewidth}\footnotesize 900点离散网格；描述性支持条件为2ΔNLL≤10，保留17点。参数限于预定等观测方差与独立过程创新族。删一范围仅使用删除后所选参数；原表附带似然元数据属于完整25期。\end{minipage}
\end{table}
